\graphicspath{{chapters/05evaluation/graphs}}

\chapter{Evaluation}

%For any practical projects, you should almost certainly have some kind
%of evaluation, and it's often useful to separate this out into its own
%chapter.

\todo{
\begin{itemize}
\item Intro to Chapter.
\item Key - compare heterogenous vs homogenous (show why heterogeneous is better)
\item Overall system performance (key graphs only).
\end{itemize}
}

% inner loop by itself
\section{Inner Loop}
\label{chap:}

We test the \textit{inner loop} using four test cases (island layouts) shown in \autoref{tab:05-gpu-island}, which include two homogeneous scenarios and two heterogeneous scenarios.
One scenario from both homogeneous and heterogeneous uses GPU islands of the same size, the other uses GPU islands of different sizes.
These islands are crafted to show the importance of selecting appropriate island sizes and potential value of heterogeneous clusters.

\begin{table}[ht]
  \centering
  \small
  \begin{tabular}{@{} l l l l @{}}
    \toprule
    Island                      & GPU Model     & Sizes (Size × Count)                & Cost         \\
    \midrule
    (a) Homogeneous Same Size       & DGX-B300      & 32 × 8                              & —            \\
    (b) Homogeneous Different Sizes & DGX-B300      & 32 × 2, 16 × 4, 8 × 8, 4 × 16        & —            \\
    \midrule
    \multirow{4}{*}{(c) Heterogeneous Same Size}
                                   & DGX-B300      & 32 × 2                              & —            \\
                                   & Rubin-NVL144  & 32 × 2                              & —            \\
                                   & H800          & 32 × 2                              & —            \\
                                   & H20           & 32 × 2                              & —            \\
    \addlinespace
    \multirow{4}{*}{(d) Heterogeneous Different Sizes}
                                   & DGX-B300      & 32 × 1, 16 × 1, 8 × 1, 4 × 2         & —            \\
                                   & Rubin-NVL144  & 32 × 1, 16 × 1, 8 × 1, 4 × 2         & —            \\
                                   & H200          & 8 × 4, 4 × 8                        & —            \\
                                   & H20           & 8 × 4, 4 × 8                        & —            \\
    \bottomrule
  \end{tabular}
  \caption{GPU island test cases for \textit{inner loop} analysis, with two homogeneous scenarios and two heterogeneous scenarios}
  \label{tab:05-gpu-island}
\end{table}

As the \textit{inner loop} solution generation is deterministic we only execute each test case once, unless stated otherwise.
This is because the linear solver will produce the exact same solution upon every execution, given the same input variables.

\clearpage
\subsection{Solver Results}
% TODO - Must also include all the input parameters the solver 

We first simultaneously select the optimum parallelism configuration and benchmark the provided GPU islands for both prefill and decode.
The resultant throughput values are displayed in \autoref{fig:05-inner-benchmark}, where we can see that every island class is able to contribute to prefill throughput, but not all islands are capable of decode.
This is due to the decode phase requiring more memory to execute, hence not all islands are compatible.
% TODO - Add explanation as to where the ranges are from.

\begin{figure}[ht]
    \centering
    \includegraphics[width=\textwidth]{inner_benchmark_all_inventories.pdf}
    \caption[Benchmark results for every island test case] 
    {Benchmark results for every island test case}
    \label{fig:05-inner-benchmark}
\end{figure}

Using the collected benchmark values, we populate the $b^{\mathrm{prefill}}_{i j}$ and $b^{\mathrm{decode}}_{i j}$ variables in the linear solver.
Then, using artificial workload as previously shown in \autoref{fig:04-sample-prefill-pdf}, we execute the solve function to find the optimum $x_{\mathrm{prefill},\,i j}$ and $x_{\mathrm{decode},\,i j}$ values.
The resultant assignment percentages are displayed in \autoref{fig:05-inner-assignment} for both the prefill and decode phase.

In the homogenous scenario, the assignment allocation percentages closely resemble the target workload.
In contrast, heterogenous assignment allocation percentages appear random.
This is to be expected, as it reflects the heterogeneous nature of island performance, where one island might be allocated few ranges due to its ``slow'' performance.
Also, for all scenarios, one large island is generally allocated for decode and many islands for prefill, this displays that prefill is the bottleneck for throughput.

\begin{figure}[ht]
    \centering
    \includegraphics[width=\textwidth]{inner_assignment_all_inventories.pdf}
    \caption[Assignment values percentage (per range) for each every island test case]
    {Assignment values percentage (per range) for each every island test case, using artificial workload as described in \autoref{fig:04-sample-prefill-pdf}}
    \label{fig:05-inner-assignment}
\end{figure}

Collating the assignment percentages with throughout measurements (\autoref{fig:05-inner-assignment} and \autoref{fig:05-inner-benchmark}) we can plot the actual throughput per island, shown in \autoref{fig:05-inner-throughput}.

\begin{figure}[ht]
    \centering
    \includegraphics[width=\textwidth]{inner_throughput_all_inventories.pdf}
    \caption[Throughput per range for each every island test case]
    {Throughput per range for each every island test case, using artificial workload as described in \autoref{fig:04-sample-prefill-pdf}}
    \label{fig:05-inner-throughput}
\end{figure}

\clearpage
\subsection{Evaluator Results}

Previous plots demonstrate the allocation results using direct outputs from the linear solver at the specified range resolution.
We need to evaluate the solutions using the complete expected input sequence length distribution $p(s \mid w)$ to highlight any differences in performance.
These results are displayed in \autoref{fig:05-inner-trace-throughput}, where \textit{evaluated} refers to results from the evaluator and \textit{model} refers to results directly from the solver.
We also include the workload \textit{trace} as a reference to what the target throughput distribution is.

\begin{figure}[ht]
    \centering
    \includegraphics[width=\textwidth]{inner_trace_throughput}
    \caption[Inner Trace Throughput]
    {Inner Trace Throughput}
    \label{fig:05-inner-trace-throughput}
\end{figure}


%display how the resolution impacts the solution

%describe the graphs, like why each graph is different 


%a plot of the throughput against cost for each island test case (show why heterogenous is better)


\subsubsection{Evaluator Method}
The true throughput is calculated by first benchmarking each exact sequence length (whilst also selecting the best parallelism configuration), for each island using the same method as described in \autoref{chap:inner-loop}.
Then, the throughputs are weighted with the $p$ value of each sequence length $s$ along with the island percentage allocation $x$, of the range the sequence length resides within.

% resolution
\subsection{Resolution Impact}

\todo{
\begin{itemize}
\item To Complete.
\item How resolution impacts solution quality vs time.
\end{itemize}
}

% resolution
\subsection{Objective Impact}

\todo{
\begin{itemize}
\item To Complete. 
\item How the objective function has been tuned.
\end{itemize}
}

% outer loop + inner loop
\clearpage
\section{Outer Loop}

\todo{
\begin{itemize}
\item Test setup.
\item Test approach (repeat 10 times).
\end{itemize}
}

% Batch Size
\subsection{Batch Size Analysis}

\todo{
\begin{itemize}
\item How batch size influences results (variance).
\item Solution achieved in 20 iterations.
\end{itemize}
}

%Explain \autoref{fig:05-batch-analysis}
\begin{figure}[ht]
    \centering
    \includegraphics[width=\textwidth]{outer_batch_errorbars.pdf}
    \caption[Batch Size Analysis]
    {Batch Size Analysis}
    \label{fig:05-batch-analysis}
\end{figure}

\todo{
\begin{itemize}
\item How long on average to best solution.
\end{itemize}
}

\begin{figure}[ht]
    \centering
    \includegraphics[width=0.5\textwidth]{outer_batch_deltas.pdf}
    \caption[Batch Size Analysis (Time)]
    {Batch Size Analysis (Time)}
    \label{fig:05-batch-deltas}
\end{figure}

% Skew
\clearpage
\subsection{Skew Analysis}

\todo{
\begin{itemize}
\item How changing the search params (skew) affects solution.
\item 
\end{itemize}
}

\begin{figure}[ht]
    \centering
    \includegraphics[width=\textwidth]{outer_skew_error_area.pdf}
    \caption[Skew Analysis]
    {Skew Analysis}
    \label{fig:05-skew-analysis}
\end{figure}

% Skew
\subsection{Number of Islands}

\todo{
\begin{itemize}
\item How changing the search params (number of islands) affects solution.
\end{itemize}
}

% Introduce where the traces have come from
\clearpage
\section{Real-World}

\todo{
\begin{itemize}
\item End-to-end analysis with real-world traces.
\end{itemize}
}

\subsection{Traces}
\graphicspath{{chapters/04design/graphs}}

% The main set of traces (Azure LLM)
\subsubsection{Azure LLM}

Source \autocite{stojkovicDynamoLLMDesigningLLM2024}

\begin{figure}[ht]
    \centering
    \includegraphics[width=0.8\textwidth]{traces_azure_context_tokens_distribution.pdf}
    \caption[Context Tokens]
    {Context Tokens}
    \label{fig:05-traces-azure-context-tokens-distribution}
\end{figure}

\begin{figure}[ht]
    \centering
    \includegraphics[width=0.8\textwidth]{traces_azure_generated_tokens_distribution.pdf}
    \caption[Decode Tokens]
    {Decode Tokens}
    \label{fig:05-traces-azure-generated-tokens-distribution}
\end{figure}

\begin{figure}[ht]
    \centering
    \begin{subfigure}[t]{0.48\textwidth}
        \centering
        \includegraphics[width=\textwidth]{traces_azure_conv_gen_vs_context.pdf}
        \caption{Converstation Trace}
        \label{fig:05-traces-azure-conv-gen-vs-context}
    \end{subfigure}
    \hfill
    \begin{subfigure}[t]{0.48\textwidth}
        \centering
        \includegraphics[width=\textwidth]{traces_azure_code_gen_vs_context.pdf}
        \caption{Code Trace}
        \label{fig:05-traces-azure-code-gen-vs-context}
    \end{subfigure}
    \caption{Comparison of Converstation and Code}
    \label{fig:05-traces-azure-gen-vs-context}
\end{figure}

% Another set of traces (BurstGPT)
\clearpage
\subsubsection{BurstGPT}

Source \autocite{wangBurstGPTRealworldWorkload2025}

\begin{figure}[ht]
    \centering
    \includegraphics[width=0.8\textwidth]{traces_burst_context_tokens_distribution.pdf}
    \caption[Context Tokens]
    {Context Tokens}
    \label{fig:05-traces-burst-context-tokens-distribution}
\end{figure}

\begin{figure}[ht]
    \centering
    \includegraphics[width=0.8\textwidth]{traces_burst_generated_tokens_distribution.pdf}
    \caption[Decode Tokens]
    {Decode Tokens}
    \label{fig:05-traces-burst-generated-tokens-distribution}
\end{figure}

\begin{figure}[ht]
    \centering
    \begin{subfigure}[t]{0.48\textwidth}
        \centering
        \includegraphics[width=\textwidth]{traces_burst_conv_gen_vs_context.pdf}
        \caption{Converstation Trace}
        \label{fig:05-traces-burst-conv-gen-vs-context}
    \end{subfigure}
    \hfill
    \begin{subfigure}[t]{0.48\textwidth}
        \centering
        \includegraphics[width=\textwidth]{traces_burst_code_gen_vs_context.pdf}
        \caption{Code Trace}
        \label{fig:05-traces-burst-code-gen-vs-context}
    \end{subfigure}
    \caption{Comparison of Converstation and Code}
    \label{fig:05-traces-burst-gen-vs-context}
\end{figure}
 
 \clearpage
 \subsection{Results}
 
 \todo{
\begin{itemize}
\item How the traces perform.
\end{itemize}
}

\clearpage
\section{Heterogeneous vs Homogenous}
 
\todo{
\begin{itemize}
\item Throughput vs Cost of Heterogeneous vs Homogenous cluster.
\end{itemize}
}
